% Generated draft: requires mathematical and source review.
\documentclass[11pt]{article}
\usepackage[margin=25mm]{geometry}
\usepackage{fontspec}
\setmainfont{Latin Modern Roman}
\usepackage{amsmath,amssymb,mathtools,unicode-math}
\setmathfont{Latin Modern Math}
\usepackage{enumitem,needspace,xurl,hyperref}
\setlength{\parindent}{0pt}
\setlength{\parskip}{5pt}
\setlength{\emergencystretch}{3em}
\newcommand{\N}{\mathbb{N}}
\newcommand{\Z}{\mathbb{Z}}
\newcommand{\Q}{\mathbb{Q}}
\newcommand{\R}{\mathbb{R}}
\newcommand{\C}{\mathbb{C}}
\newcommand{\F}{\mathbb{F}}
\newcommand{\A}{\mathbb{A}}
\newcommand{\PP}{\mathbb{P}}
\newcommand{\E}{\mathbb{E}}
\newcommand{\eps}{\varepsilon}
\newcommand{\dd}{\,\mathrm{d}}
\begin{document}
\section*{UnsolvedMath ID: 801103}
\textbf{Source problem code:} catalogue entry
\subsection{Definitions and mathematical statement}
For a discrete group \(G\), let \(\lambda_G\) be its left regular representation on \(\ell^2(G)\), given by \(\lambda_G(g)\delta_h=\delta_{gh}\). Its group von Neumann algebra is
\[
L(G)=\{\lambda_G(g):g\in G\}'',
\]
with canonical trace \(\tau_G(x)=\langle x\delta_e,\delta_e\rangle\). An infinite group \(G\) is \emph{icc} if every conjugacy class other than \(\{e\}\) is infinite; in this case \(L(G)\) is a type \(\mathrm{II}_1\) factor. For an icc group, inner amenability means the existence of a conjugation-invariant finitely additive probability measure on all subsets of \(G\setminus\{e\}\). Non-inner amenability means that no such measure exists.

The supplied record does not define \emph{pseudocompact}. Write \(\mathcal P(M)\) solely as a placeholder for the source's intended assertion that the von Neumann algebra \(M\) is pseudocompact. This notation does not resolve the missing definition. With this qualification, the two questions attributed to Elayavalli are:
\begin{enumerate}
\item Does there exist a group \(G\) such that \(\mathcal P(L(G))\)?
\item Does there exist an infinite icc group \(G\) that is not inner amenable and satisfies \(\neg\mathcal P(L(G))\)?
\end{enumerate}
The group von Neumann algebra notation is interpreted here in its usual discrete-group setting. The record supplies no countability hypothesis. Although the source title concerns type \(\mathrm{II}_1\) factors, the first question does not explicitly require \(G\) to be icc or \(L(G)\) to be a type \(\mathrm{II}_1\) factor; this intended restriction remains unresolved. Consequently, the missing definition and domain qualifications must be checked before these questions become fully specified mathematical statements.

The registry labels the overall problem open, while the supplied statement separately reports that the second question was solved by J. Peterson, citing \texttt{peterson2026ultraproductapproximationspropertyt}. These are supplied status annotations, not independently verified conclusions; the record does not explicitly state the answer established for the second question.
\subsection{Short English statement}
The source asks whether a group von Neumann algebra can be pseudocompact, and whether an icc group that is not inner amenable can have a group von Neumann algebra that is not pseudocompact. The supplied record reports the second question solved by J. Peterson. The meaning of pseudocompact and the intended restrictions on groups in the first question require clarification.
\subsection{Sources}
\begin{itemize}
\item Local UnsolvedMath registry: permanent ID \(801103\), source code \texttt{catalogue entry}, title \emph{Ultraproduct II\_1 Factors --- Problem 4}; set: AMR Open Problem Lists. Attribution supplied by the registry: Elayavalli, proposed in 2025. The overall status is recorded as open.
\item Supplied source list: Elayavalli, \emph{50 Open Problems: Ultraproduct II\_1 Factors} (2025), Problem 4. Source URL: \url{https://arxiv.org/abs/2511.20377}. The record reports access on 2026-10-16 and extraction by \texttt{arxiv\_latex}; this provenance is supplied metadata, not an independently verified visit or literature assessment.
\item Supplied resolution attribution for the second question: J. Peterson, reference key \texttt{peterson2026ultraproductapproximationspropertyt}. No full bibliographic entry, URL, theorem statement, or explicit answer for this reference is supplied.
\end{itemize}
\end{document}
